% Vietnamese translation for OI Public Library
% Source-PDF: IOI中国国家候选队论文/2002/金恺.pdf
\documentclass[11pt]{article}
\usepackage{oipl-vn}

\title{Bàn về ứng dụng của thuật toán luồng mạng}
\author{Jin Kai}
\date{}

\begin{document}
\maketitle

\origtitle{A discussion of applications of network-flow algorithms}
\sourcepath{IOI China candidate team papers / 2002 / Jin Kai.pdf}

\paragraph{Từ khóa.}
Luồng mạng; xây dựng mô hình; tối ưu hóa.

\tableofcontents

\section{Dẫn nhập}

Thuật toán đồ thị giữ vai trò rất quan trọng trong các kỳ thi tin học. Số
nhánh của nó rất nhiều, phạm vi ứng dụng rất rộng, đến mức các nhóm thuật
toán khác khó sánh kịp. Thuật toán luồng mạng là một nhánh quan trọng trong
thuật toán đồ thị; nó có đặc điểm rõ rệt và tác dụng nổi bật, nên phạm vi ứng
dụng rất rộng. Trong các kỳ thi tin học ở nhiều cấp những năm gần đây, các
bài toán dùng luồng mạng xuất hiện liên tiếp, và vị trí của nó sẽ còn ngày
càng quan trọng hơn.

Mục đích của bài viết này là thông qua việc phân tích một số ví dụ ứng dụng,
từng bước trình bày cách xây dựng mô hình, nguyên tắc tối ưu hóa và phương
pháp tối ưu hóa cho thuật toán luồng mạng, từ đó giúp ta hiểu thuật toán này
sâu hơn và triệt để hơn.

Ở đây giả thiết người đọc đã quen thuộc với khung cơ bản và cách giải các
bài toán luồng mạng. Nội dung sau chủ yếu ghi lại một số kinh nghiệm trong
quá trình giải bài, tập trung vào phạm vi ứng dụng, cách áp dụng cụ thể và
nhiều kỹ thuật tối ưu hóa thường gặp.

\section{Ví dụ 1: bài toán đua xe}

\subsection{Phát biểu bài toán}

P và Q đều là cao thủ xe đua bốn bánh. Mỗi người có \(N\) chiếc xe. Để phân
cao thấp, họ hẹn nhau tổ chức một cuộc thi.

Cuộc thi có \(M\) chặng. Người thắng nhiều chặng hơn sẽ giành chức vô địch
chung cuộc. Trong mỗi chặng, mỗi người chọn một chiếc xe của mình để thi đấu.
Thắng thua của một chặng phần lớn phụ thuộc vào tính năng của hai xe, nhưng
vì nhiều nguyên nhân, chẳng hạn sự cản trở giữa hai xe, tính năng không phải
là chỉ tiêu quyết định duy nhất. Có thể xuất hiện tình huống \(A\) thắng
\(B\), \(B\) thắng \(C\), \(C\) thắng \(D\), nhưng \(D\) lại thắng \(A\).
May mắn là có một máy thông minh: chỉ cần đưa vào hai chiếc xe, nó lập tức
phán đoán được xe nào thắng. Trước cuộc thi, máy đã kiểm tra từng cặp gồm một
xe của P và một xe của Q, rồi đưa bảng thắng thua vào máy tính.

Ngoài ra, do hao mòn khi thi đấu, mỗi chiếc xe có tuổi thọ riêng, tức số
chặng mà nó có thể tham gia. Tuổi thọ của các xe có thể khác nhau, nhưng
không vượt quá \(50\) chặng. Hai chiếc xe bất kỳ nhiều nhất chỉ được đối đầu
một lần.

Cho \(N,M\), với \(1\le N\le 100\), \(1\le M\le 1000\), một bảng thắng thua
\(N\times N\), tuổi thọ của từng chiếc xe, và giả thiết rằng ở mỗi chặng cả
hai người đều còn xe có thể chọn. Hãy tính số chặng nhiều nhất mà P có thể
thắng.

\subsection{Xây dựng mạng}

Ta xây dựng một mạng như sau.

\begin{enumerate}
  \item Tạo \(N\) đỉnh biểu diễn \(N\) chiếc xe của P, đánh số từ \(1\) đến
  \(N\).
  \item Tạo \(N\) đỉnh biểu diễn \(N\) chiếc xe của Q, đánh số từ \(N+1\)
  đến \(2N\).
  \item Nếu xe thứ \(i\) của P có thể thắng xe thứ \(j\) của Q, thêm cung
  \(i\to N+j\) có dung lượng \(1\). Dung lượng này biểu diễn điều kiện hai
  chiếc xe nhiều nhất chỉ gặp nhau một lần.
  \item Thêm nguồn \(S\). Với mỗi \(i\in\{1,\ldots,N\}\), nối
  \(S\to i\), dung lượng bằng tuổi thọ của xe thứ \(i\) của P.
  \item Thêm đích \(T\). Với mỗi \(j\in\{1,\ldots,N\}\), nối
  \(N+j\to T\), dung lượng bằng tuổi thọ của xe thứ \(j\) của Q.
  \item Thêm một nguồn mới \(S_2\), nối \(S_2\to S\) với dung lượng \(M\),
  biểu diễn tổng cộng nhiều nhất \(M\) chặng.
\end{enumerate}

Khi đó mỗi đơn vị luồng trên một đường
\[
  S_2\to S\to i\to N+j\to T
\]
tương ứng với việc chọn xe \(i\) của P thắng xe \(j\) của Q trong một chặng.
Giá trị luồng cực đại chính là số chặng thắng nhiều nhất của P.

Có thể xem cấu trúc mạng gồm bốn loại cung:
\[
\begin{array}{c|l|l}
\text{Loại} & \text{Cung} & \text{Dung lượng}\\ \hline
1 & S_2\to S & M\\
2 & S\to i,\quad 1\le i\le N & \text{tuổi thọ xe }i\text{ của P}\\
3 & i\to N+j & 1,\text{ nếu P}_i\text{ thắng Q}_j\\
4 & N+j\to T,\quad 1\le j\le N & \text{tuổi thọ xe }j\text{ của Q}
\end{array}
\]

\subsection{Tối ưu hóa}

Nếu dùng cách thông thường, chỉ riêng việc lưu cả lưu lượng và dung lượng đã
cần \((2N+3)(2N+3)\cdot 2\) bit. Nhưng mạng ở đây chỉ có bốn loại cung như
trên; số cung nhiều nhất chỉ là
\[
  1+N+N+N^2=(N+1)^2.
\]
Vì vậy ta không cần lưu cả ma trận đầy đủ.

Hơn nữa, nguồn \(S_2\) và cung \(S_2\to S\) cũng có thể bỏ. Chỉ cần quy định
thuật toán tăng luồng nhiều nhất \(M\) lần; điều này tương đương với giới hạn
tổng số chặng.

\section{Ví dụ 2: điều độ tàu hỏa}

\subsection{Phát biểu bài toán}

Một ga hàng hóa có \(n\) đường ray, \(n\le 20\). Do chiều dài đường ray có
hạn, tại một thời điểm mỗi đường ray nhiều nhất chỉ đỗ được một đoàn tàu
hàng. Sau khi ga vận hành bình thường, mỗi ngày có \(m\) đoàn tàu hàng đi qua
ga, \(m\le 100\). Đoàn tàu thứ \(i\) đến ga tại thời điểm \(\mathrm{Reach}[i]\)
và chở hàng có giá trị \(\mathrm{Cost}[i]\).

Nếu cho phép tàu \(i\) vào ga, ga BackStreet nhận được lợi nhuận
\(1\%\cdot \mathrm{Cost}[i]\). Nhưng do cần thời gian bốc dỡ, tàu này sẽ dừng
ở ga trong \(\mathrm{Stay}[i]\) đơn vị thời gian, và trong khoảng thời gian đó
nó chiếm một đường ray. Nếu không cho vào ga, tàu đi thẳng ra ngoài và ga
không thu được gì.

Nhiệm vụ là sắp xếp hợp lý việc tàu vào ga hoặc ra thẳng sao cho tổng lợi
nhuận của ga lớn nhất. Nếu tàu \(a\) vừa rời đường ray thứ \(i\) đúng lúc
tàu \(b\) đến, tức
\[
  \mathrm{Reach}[a]+\mathrm{Stay}[a]=\mathrm{Reach}[b],
\]
thì tàu \(b\) không được vào đường ray thứ \(i\). Do đó một tàu có thể đứng
sau tàu \(i\) trên cùng một đường ray chỉ khi
\[
  \mathrm{Reach}[i]+\mathrm{Stay}[i]<\mathrm{Reach}[j].
\]

\paragraph{Ví dụ.}
Với \(2\) đường ray và \(5\) đoàn tàu, dữ liệu sau cho lợi nhuận lớn nhất
\(0.18\). Một cách đạt được là đưa tàu \(1,3\) vào đường ray thứ nhất và tàu
\(2,5\) vào đường ray thứ hai.

\[
\begin{array}{c|ccccc}
\text{Số hiệu tàu} & 1 & 2 & 3 & 4 & 5\\ \hline
\text{Thời điểm vào ga} & 2 & 3 & 5 & 2 & 6\\
\text{Giá trị hàng} & 5 & 4 & 6 & 2 & 3\\
\text{Thời gian dừng} & 2 & 2 & 3 & 2 & 2
\end{array}
\]

\subsection{Xây dựng mạng chi phí}

Bài toán là bài toán luồng cực đại chi phí lớn nhất với dung lượng đơn vị.
Xây dựng mạng như sau.

\begin{enumerate}
  \item Tách mỗi đoàn tàu thành hai đỉnh. Tàu \(i\) được tách thành \(i\) và
  \(i'\). Nối \(i\to i'\) với dung lượng \(1\), chi phí
  \(1\%\cdot \mathrm{Cost}[i]\). Nếu cung này có lưu lượng \(1\), nghĩa là
  tàu \(i\) được cho vào ga và sinh lợi nhuận tương ứng.
  \item Thêm nguồn \(S\). Với mỗi đỉnh \(i\), nối \(S\to i\), dung lượng
  \(1\), chi phí \(0\). Nếu cung này có lưu lượng \(1\), tàu \(i\) là đoàn
  tàu đầu tiên được đưa vào một đường ray nào đó.
  \item Thêm nguồn phụ \(S'\). Nối \(S'\to S\) với dung lượng \(n\), biểu
  diễn có \(n\) đường ray.
  \item Thêm đích \(T\). Với mỗi đỉnh \(i'\), nối \(i'\to T\), dung lượng
  \(1\), chi phí \(0\). Nếu cung này có lưu lượng \(1\), tàu \(i\) là đoàn
  tàu cuối cùng được đưa vào một đường ray nào đó.
  \item Với mọi \(i\ne j\), nếu
  \[
    \mathrm{Reach}[i]+\mathrm{Stay}[i]<\mathrm{Reach}[j],
  \]
  nối \(i'\to j\) với dung lượng \(1\), chi phí \(0\). Cung này biểu diễn
  việc sau khi tàu \(i\) rời ga, tàu \(j\) có thể dùng tiếp cùng đường ray.
\end{enumerate}

Một đường đi từ \(S'\) đến \(T\) tương ứng với chuỗi tàu dùng chung một đường
ray. Tổng số đường đi được chọn không vượt quá \(n\), và tổng chi phí của
luồng là tổng lợi nhuận của các đoàn tàu được đưa vào ga.

\subsection{Tối ưu hóa}

Thực ra không nhất thiết phải thêm nguồn phụ \(S'\) và cung \(S'\to S\). Mỗi
lần sửa lưu lượng trên đường tăng, lượng sửa đều bằng \(1\). Vì vậy chỉ cần
quy định tìm nhiều nhất \(n\) đường tăng là bảo đảm số đường ray bị chiếm
không vượt quá \(n\).

Sau tối ưu hóa thứ nhất, mọi cung đều có dung lượng \(1\), rất tiện xử lý.
Có thể dùng một biến kiểu byte để lưu cả dung lượng và lưu lượng của mỗi
cung:
\[
\begin{array}{c|l}
0 & \text{dung lượng }0,\ \text{lưu lượng }0,\\
1 & \text{dung lượng }1,\ \text{lưu lượng }0,\\
2 & \text{dung lượng }1,\ \text{lưu lượng }1.
\end{array}
\]

\section{Ví dụ 3: bài toán nhà hàng}

\subsection{Phát biểu bài toán}

Trong \(n\) ngày liên tiếp, mỗi ngày công ty có một nhu cầu khăn nhất định.
Ngày thứ \(i\) cần \(A_i\) chiếc khăn. Khăn phải được khử trùng trước mỗi lần
sử dụng; khăn mới mua đã được khử trùng.

Có hai cách khử trùng. Cách \(A\) cần \(a\) ngày, cách \(B\) cần \(b\) ngày,
với \(b>a\). Chi phí của hai cách lần lượt là \(f_a\) và \(f_b\). Giá mua một
chiếc khăn mới là \(f\), với \(f>f_a>f_b\). Hãy tìm chi phí nhỏ nhất để đáp
ứng nhu cầu mỗi ngày.

\subsection{Cách xây dựng thứ nhất}

Theo cách thông thường, ta tách mỗi ngày thành hai đỉnh: ngày thứ \(i\) gồm
đỉnh \(i\) và đỉnh \(i'\), lần lượt xem như thời điểm bắt đầu và kết thúc
ngày đó. Ngoài ra thêm nguồn \(S\) và đích \(T\).

Các cung được nối như sau.

\begin{enumerate}
  \item Nối nguồn \(S\) đến đỉnh \(1\), dung lượng từ \(0\) đến
  \(\infty\), chi phí \(f\). Lưu lượng trên cung này biểu diễn tổng số khăn
  mới được mua. Vì mua khăn mới vào bất kỳ ngày nào cũng tương đương với mua
  ngay từ ngày đầu, ta có thể mua toàn bộ khăn mới cần dùng ở ngày thứ nhất.
  \item Với \(i\le n-1\), nối \(i\to i+1\), dung lượng từ \(0\) đến
  \(\infty\), chi phí \(0\). Cung này biểu diễn khăn mới hoặc khăn đã khử
  trùng có ở đầu ngày \(i\) có thể để lại cho ngày kế tiếp.
  \item Nối \(i\to i'\). Cung này có cận dưới \(A_i\), cận trên \(A_i\), chi
  phí \(0\). Cận dưới \(A_i\) biểu diễn ngày đó ít nhất cần \(A_i\) khăn; vì
  dùng quá \(A_i\) khăn là không cần thiết, cận trên cũng đặt là \(A_i\).
  \item Nối \(i'\to T\), dung lượng từ \(0\) đến \(\infty\), chi phí \(0\).
  Cung này biểu diễn khăn đã dùng trong ngày \(i\) có thể không đem khử trùng
  nữa mà vứt bỏ trực tiếp.
  \item Với \(i<n-a\), nối \(i'\to i+a+1\), dung lượng từ \(0\) đến
  \(\infty\), chi phí \(f_a\). Cung này biểu diễn khăn đã dùng ở ngày \(i\)
  được khử trùng theo cách \(A\), và sau \(a\) ngày có thể dùng lại.
  \item Với \(i<n-b\), nối \(i'\to i+b+1\), dung lượng từ \(0\) đến
  \(\infty\), chi phí \(f_b\). Cung này biểu diễn khăn đã dùng ở ngày \(i\)
  được khử trùng theo cách \(B\), và sau \(b\) ngày có thể dùng lại.
\end{enumerate}

Sơ đồ của mạng có thể hiểu theo ba tầng ý nghĩa: cung mua khăn đi từ cửa
hàng đến đầu ngày \(1\); các cung giữ khăn nối đầu ngày này sang đầu ngày
kế tiếp; các cung khử trùng nối cuối ngày \(i\) đến đầu ngày \(i+a+1\) hoặc
\(i+b+1\); các cung bỏ khăn nối cuối mỗi ngày đến đích.

\subsection{Tối ưu hóa}

\paragraph{1. Xây dựng luồng ban đầu.}
Mạng trên là mạng có cận dưới và cận trên trên dung lượng, nên luồng ban đầu
không thể đặt là luồng \(0\). Nếu dùng cách trong sách giáo khoa về đồ thị,
tức xây thêm mạng phụ để tìm luồng khả thi ban đầu, ta phải viết thêm nhiều
mã, tốn thời gian lập trình, đồng thời không gian và thời gian chạy đều bị
lãng phí. Ở đây nên tùy bài mà xử lý.

Ý tưởng cốt lõi của luồng ban đầu là: không khử trùng, mọi nhu cầu khăn đều
được đáp ứng bằng cách mua mới. Đánh số đỉnh cuối ngày \(i\) là \(i+n\). Khi
đó có thể đặt:
\[
\begin{aligned}
  f(i,i+n)&=A_i, && 1\le i\le n,\\
  f(n-1,n)&=A_n,\\
  f(i,i+1)&=f(i+1,i+2)+f(i+1,i+n+1), && 1\le i\le n-2,\\
  f(S,1)&=f(1,2)+f(1,n+1),\\
  f(i+n,T)&=A_i, && 1\le i\le n.
\end{aligned}
\]
Các cung khử trùng ban đầu có thể đặt lưu lượng \(0\). Luồng này đúng với
kịch bản mua đủ khăn mới và sau khi dùng thì vứt bỏ.

\paragraph{2. Lưu trữ không gian.}
Trước đây, các thuật toán luồng mạng thường dùng ma trận kề hoặc danh sách
kề để lưu mạng. Với bài này, \(n\) có thể đến \(1000\), nên chắc chắn không
thể dùng ma trận kề. Nếu dùng danh sách kề, độ phức tạp lập trình lại tăng.
Quan sát cấu trúc mạng, mọi cung chỉ thuộc năm loại:
\[
\begin{array}{c|l}
1 & i-1\to i,\quad 1\le i\le n;\quad 0\to 1\text{ biểu diễn }S\to 1,\\
2 & i\to i+n,\quad 1\le i\le n,\\
3 & i+n\to T,\quad 1\le i\le n,\\
4 & i+n\to i+a+1,\quad 1\le i<n-a,\\
5 & i+n\to i+b+1,\quad 1\le i<n-b.
\end{array}
\]
Dung lượng, chi phí, và lưu lượng của cung loại \(2\) đều cố định. Vì vậy chỉ
cần dùng bốn mảng kích thước \(1..n\) để lưu lưu lượng của bốn loại cung còn
lại. Nếu mỗi giá trị dùng \(2\) byte, không gian chỉ khoảng
\[
  \frac{4\cdot 1000\cdot 2}{1024}\approx 8\text{ KB}.
\]

\section{Ví dụ 4: mạng tình báo cuối cùng}

\subsection{Phát biểu bài toán}

Bài này được dẫn chiếu là bài thứ nhất của vòng thi CTSC 2001, tên
\textit{Agent}. Các slide không chép lại toàn bộ đề, mà tập trung vào mô
hình luồng mạng và cách chuyển hàm mục tiêu dạng tích thành bài toán chi phí
luồng.

\subsection{Xây dựng mạng}

Các ký hiệu dùng trong mô hình gồm:
\begin{itemize}
  \item \(K\): giới hạn tổng thể từ tổng bộ đồng minh;
  \item \(AM_i\): lượng liên hệ mà tổng bộ đồng minh có thể thực hiện với
  điệp viên \(i\);
  \item \(AS_i\): hệ số an toàn của liên hệ từ tổng bộ đồng minh đến điệp
  viên \(i\);
  \item \(M_{i,j}\): số lượng thông tin mà điệp viên \(i\) có thể truyền cho
  điệp viên \(j\);
  \item \(S_{i,j}\): hệ số an toàn của cung truyền tin từ \(i\) đến \(j\).
\end{itemize}

Gọi \(C_{u,v}\) là dung lượng và \(W_{u,v}\) là trọng số an toàn trên cung
\((u,v)\). Các bước chính để lập mạng là:

\begin{enumerate}
  \item Thêm nguồn \(p\).
  \item Thêm đỉnh phụ \(p'\), đặt
  \[
    C_{p,p'}=K,\qquad W_{p,p'}=1.
  \]
  Nếu tổng bộ đồng minh có thể liên hệ với điệp viên \(i\), tức \(AM_i>0\),
  thì đặt
  \[
    C_{p',i}=AM_i,\qquad W_{p',i}=AS_i,
    \qquad i=1,2,\ldots,n.
  \]
  \item Nếu điệp viên \(i\) có thể truyền tin cho điệp viên \(j\), tức
  \(M_{i,j}>0\), thì đặt
  \[
    C_{i,j}=M_{i,j},\qquad W_{i,j}=S_{i,j},
    \qquad i,j=1,2,\ldots,n.
  \]
  \item Thêm đích \(q\). Nếu điệp viên \(i\) có thể liên hệ với bộ phận tình
  báo của đối phương, thì đặt
  \[
    C_{i,q}=K,\qquad W_{i,q}=1,
    \qquad i=1,2,\ldots,n.
  \]
\end{enumerate}

\subsection{Phân tích hàm mục tiêu}

Bài toán luồng cực đại chi phí nhỏ nhất tìm cách làm nhỏ nhất tổng
\[
  \sum_{(u,v)} \text{chi phí}(u,v)\cdot \text{lưu lượng}(u,v).
\]
Trong bài này, mục tiêu lại là làm lớn nhất tích các lũy thừa của hệ số an
toàn theo lưu lượng. Nếu \(F_{i,j}\) là lượng thông tin thực sự truyền từ
\(i\) đến \(j\), thì một dạng mục tiêu của bài là
\[
  \max\left\{
    \prod_{i=1}^{n}\prod_{j=1}^{n} S_{i,j}^{F_{i,j}}
    \cdot
    \prod_{i=1}^{n} AS_i^{F_{p',i}}
  \right\}.
\]
Lấy đối số logarit và đổi dấu, mục tiêu trên chuyển thành
\[
  \min\left\{
    \sum_{i=1}^{n}\sum_{j=1}^{n}
      \bigl(-F_{i,j}\ln S_{i,j}\bigr)
    +
    \sum_{i=1}^{n}
      \bigl(-F_{p',i}\ln AS_i\bigr)
  \right\}.
\]
Từ đây có thể thấy bài toán đã chuyển từ ``tìm tích'' thành ``tìm tổng'', phù
hợp với mô hình luồng cực đại chi phí nhỏ nhất. Vì thế ta có thể dùng thuật
toán chi phí luồng để giải.

Cụ thể, với mỗi cung \((a,b)\) trong mạng, đặt chi phí mới
\[
  W'_{a,b}=-\ln W_{a,b}.
\]
Mặt khác, chỉ cần \(M_{i,j}>0\) thì \(0<S_{i,j}\le 1\), do đó
\(\ln S_{i,j}\le 0\) và \(W'_{i,j}\ge 0\). Vì vậy khi tìm đường tăng có chi
phí nhỏ nhất, ta vẫn có thể dùng các phương pháp thường dùng cho bài toán
chi phí luồng.

Nếu \(F\) là luồng cuối cùng và \(C\) là dung lượng, mục tiêu cuối có thể
viết dưới dạng
\[
  \prod_{i=1}^{n}\prod_{j=1}^{n}
    e^{W_{i,j}(F_{i,j}-C_{i,j})}
  \cdot
  \prod_{i=1}^{n}
    e^{W_{p',i}(F_{p',i}-C_{p',i})}.
\]
Công thức này tương ứng với việc khôi phục giá trị mục tiêu sau khi đã giải
bài toán trên thang logarit.

\section{Kết luận}

So với các thuật toán đồ thị khác, mô hình luồng mạng phức tạp hơn, độ phức
tạp lập trình cao hơn, và hệ số hằng của thuật toán cũng lớn hơn.

Ưu điểm của mô hình luồng mạng là nó chứa được rất nhiều yếu tố, đặc biệt là
nhiều loại trọng số: không chỉ có trọng số biểu diễn dung lượng và lưu lượng,
mà còn có trọng số biểu diễn chi phí; dung lượng không chỉ có cận trên mà còn
có thể có cận dưới. Chính các yếu tố biến hóa này làm cho mô hình luồng mạng
có phạm vi ứng dụng rất rộng trong việc giải quyết các bài toán thực tế và
nghiên cứu bài toán đồ thị. Nó thường giải quyết tốt một số bài toán mà tìm
kiếm hoặc quy hoạch động khó xử lý, trông giống như các bài toán NP.

Khó khăn lớn nhất khi áp dụng luồng mạng vào một bài cụ thể là xây dựng mô
hình; khó khăn tiếp theo là tối ưu hóa thuật toán.

Việc xây dựng mô hình không có khuôn mẫu sẵn. Ta phải phân tích theo điều
kiện cụ thể của đề bài. Điều này đòi hỏi hiểu rõ các tính chất của nhiều dạng
luồng mạng, biết đúc kết kinh nghiệm, và phát huy tính sáng tạo. Nói chung,
phương pháp được dùng nhiều nhất là tách đỉnh.

Tối ưu hóa là một khâu quan trọng của thuật toán. Nó không thể tiến bộ trong
một sớm một chiều, mà phải dựa vào tích lũy kinh nghiệm.

\end{document}
